\section{TailSpline: Construction and Frozen Evaluation}
\label{sec:tailspline-details}
\subsection{Finite-grid proof}
Let $D\epsilon=(\epsilon_{q+1}-\epsilon_q)_{q=1}^{n-1}$ and $H=D^\top D+e_ne_n^\top$. A zero value of $\epsilon^\top H\epsilon$ requires constant increments and a zero terminal increment, hence $\epsilon=0$. Thus $H$ is positive definite. For $n>1$, the proposed $\epsilon_q=c(n+q)(n-q+1)$, with $c=3/[n(n+1)(2n+1)]$, satisfies
\[
\epsilon_1-\epsilon_2=2c,\quad
2\epsilon_q-\epsilon_{q-1}-\epsilon_{q+1}=2c,\quad
2\epsilon_n-\epsilon_{n-1}=2c.
\]
Therefore $H\epsilon=2c\mathbf1$, the equality-constrained stationarity equation. The increments are positive and sum to one; summing them yields Eq.~\eqref{eq:tailspline}. The objective minimum is $2c=6/[n(n+1)(2n+1)]$. For $n=1$, the constraint gives $\epsilon_1=1$, agreeing with the formulas. Positivity makes additional nonnegativity constraints inactive.

Adding $\epsilon_1^2$ yields the symmetric objective, with unique solution
\begin{equation}
\epsilon_q^{\rm BM}=\frac{6q(n-q+1)}{n(n+1)(n+2)},\qquad
m_q^{\rm BM}=\frac{q(q+1)(3n+2-2q)}{n(n+1)(n+2)}.
\label{eq:bm-formula}
\end{equation}
Both profiles are monotone with $m_0=0,m_n=1$. The portable CPU check independently solves the constrained tridiagonal system and verifies each closed form.

\paragraph{What the objective smooths.}
On the geometric native grid, Eq.~\eqref{eq:increment-log-gap} makes each change between neighboring log gaps equal to $(\epsilon_{q+1}-\epsilon_q)\log s$. The terminal change is $-\epsilon_n\log s$ because the fully interpolated tail has constant displacement. Thus $(\log s)^2$ times Eq.~\eqref{eq:tailspline-objective} is exactly the squared log-gap variation within the transition and at its low-frequency junction. The high-frequency junction is deliberately unpenalized.

At the terminal gap,
\[
\epsilon_n^{\rm Pro}=\frac2{n+1},\quad
\epsilon_n^{\rm TS}=\frac6{(n+1)(2n+1)},\quad
\frac{\epsilon_n^{\rm TS}}{\epsilon_n^{\rm Pro}}=\frac3{2n+1}.
\]
The opposite entry behavior is also explicit: $\epsilon_1^{\rm Pro}=2/[n(n+1)]$, whereas $\epsilon_1^{\rm TS}=3/(2n+1)$. TailSpline trades a larger entry jump for a smaller terminal jump. A useful boundary prior need not minimize both. In the continuum limit, $m^{\rm TS}(u)=(3u-u^3)/2$ solves $\min\int_0^1(m''(u))^2du$ under $m(0)=0,m(1)=1,m'(1)=0$; the natural entry condition is $m''(0)=0$. The finite-grid formula is used in every model experiment, not this limiting polynomial.

\subsection{An exact equal-displacement control}
\label{sec:tailspline-dose-control}
Use local indices $q=0,\ldots,n$. Write $T_q=m_q^{\rm TS}$, $B_q=m_q^{\rm BM}$, $U_q=q/n$, and $P_q=q(q+1)/[n(n+1)]$. Define the front-loaded profile $F_q=2U_q-P_q$ and the control
\begin{equation}
C_q=(1-w_n)U_q+w_nF_q,\qquad w_n=\frac{3n}{2(2n+1)}.
\label{eq:tailspline-dose-control}
\end{equation}
Substitution of the finite-grid formulas gives $T=(1-w_n)B+w_nF$. Subtracting $C$ therefore yields Eq.~\eqref{eq:tailspline-shape-contrast}, with the explicit difference
\begin{equation}
T_q-C_q=(1-w_n)(B_q-U_q)=\frac{q(n-q)(2q-n)}{2n(n+1)(2n+1)}.
\label{eq:tailspline-antisymmetric-contrast}
\end{equation}
This cumulative-profile difference is antisymmetric under $q\mapsto n-q$, vanishes at both endpoints, and sums to zero. Since $U$ and $F$ have positive increments and $0<w_n<1$, $C$ is monotone with $C_0=0,C_n=1$. Extending both $T$ and $C$ by the same constant outer bands gives identical endpoints and total log displacement at a common scale. Their equal increment centroids follow from Eq.~\eqref{eq:displacement-centroid}, rather than imposing a further constraint. For $n=1,2$, $B=U$ and $T=C$; for wider bands the contrast tests higher-order shape.

The reported TailSpline--MrPro results below change total displacement and shape together. A matched $T$--$C$ task comparison would isolate the zero-total correction above; the algebra alone supplies no task-score ordering and does not identify the tail jump as its sole mediator. No $T$--$C$ task result is included here. The independent CPU check in \path{figs/allocation_design.py} verifies the finite-sum and profile identities with exact rational arithmetic.

\subsection{Protocol and complete length responses}
Frozen Llama-$3$-$8$B-Instruct has native window $8192$. Both arms use exact finite-grid profiles, $s=4$, canonical $32/1$-turn band $[18,35]$, and cosine/sine gain $1+0.1\ln4$. Each table is installed before prefill and shared across layers, requests and evaluation lengths. Construction uses no weights, activations or calibration outputs.

The depth-balanced RULER-13 panel contains $10$ prompts per task at $8/16/32$K: $390$ paired inputs per arm. NIAH and passkey are subsets of this panel. PPL uses $32$ Proof-Pile-test and $14$ PG-19-test documents at each length: $138$ rows per arm \citep{azerbayev2022proofpile,rae2019compressive}. Both compared methods retain their $s=4$ table even at the native input length.

Normalized trapezoidal log-length AUC weights the three endpoints by $1/4,1/2,1/4$. For PPL it integrates PPL over log length, not log PPL. At each length, PPL is the exponential of total whole-sequence loss divided by total scored tokens across the 46 documents. The paired document bootstrap resamples document identities jointly across lengths, recomputing token-weighted NLL, PPL and AUC for each draw. Task aggregates weight tasks equally. The recorded paired intervals characterize the tested inputs and task families for this frozen checkpoint.
\begin{table}[ht]
\centering\small
\caption{Recorded Llama responses. NIAH uses recorded paired differences; absolute scores are not reconstructed.}
\begin{tabular}{@{}lrrr@{}}
\toprule
Endpoint & $8$K & $16$K & $32$K \\
\midrule
RULER-13 MrPro (\%) & 87.5000 & 80.0897 & 54.7308 \\
RULER-13 TailSpline (\%) & 89.7436 & 84.0513 & 57.3590 \\
NIAH TailSpline--MrPro (pp) & $-0.3125$ & $+5.9375$ & $+3.4375$ \\
PPL MrPro & 5.286429 & 4.518111 & 4.161831 \\
PPL TailSpline & 5.295283 & 4.530318 & 4.110622 \\
\bottomrule
\end{tabular}
\end{table}
\FloatBarrier
Family AUC differences are retrieval $+3.750$ points, tracking $+6.500$, aggregation approximately $+3.792$, and QA $-1.250$. Multikey-2 improves by $22.5$ points; QA-1 and multikey-1 each decrease by $5.0$. Out of $390$ inputs, TailSpline and MrPro respectively have $30$ and $37$ empty outputs, $388$ and $384$ EOS-terminated outputs, and $2$ and $6$ budget-limited outputs. Task scores include partial credit.

The portable input \path{figs/allocation_value_inputs.json} records the aggregates and source identities. The packaged constructor computes the boundaries and FP32 table from the public native grid, reference length and scale. TailSpline and MrPro share their outer bands, endpoints and gain, so their matched comparison tests interior allocation. Their total log-frequency displacement also differs; at fixed support this is a statistic of $z$. The equal-total controls in Appendix~\ref{sec:four-method-control} further isolate shape at fixed displacement.
\FloatBarrier

\subsection{Prospective OLMo confirmation}
The OLMo-$2$-0425-$1$B-Instruct comparison uses original window $4096$, canonical band $[14,32]$, $s=4$ and the same gain as Llama. Generation batch size is four. The methods and data were frozen before reading outputs from this experiment; historical BM/front-loaded work on OLMo had informed the direction. Each model has its own tokenized panel and within-model paired comparison.
\begin{table}[ht]
\centering\small
\caption{Complete OLMo length responses under the frozen two-arm contract.}
\begin{tabular}{@{}lrrr@{}}
\toprule
Endpoint & $4$K & $8$K & $16$K \\
\midrule
RULER-13 MrPro (\%) & 32.5128 & 14.8462 & 7.2436 \\
RULER-13 TailSpline (\%) & 72.6667 & 72.5513 & 48.6154 \\
NIAH MrPro (\%) & 33.1250 & 13.1250 & 7.1875 \\
NIAH TailSpline (\%) & 88.7500 & 87.8125 & 60.0000 \\
PPL MrPro & 13.795244 & 14.356459 & 12.827048 \\
PPL TailSpline & 11.152472 & 10.117407 & 8.536196 \\
\bottomrule
\end{tabular}
\end{table}
\FloatBarrier
Full-13 AUC difference is $0.492340$ with paired $95\%$ interval $[0.453365,0.530449]$; NIAH difference $0.644531$ has interval $[0.592188,0.694531]$; PPL difference $-3.852932$ has interval $[-4.775042,-3.125459]$. Passkey AUC is $0.950$ versus $0.375$. All $13$ task differences are positive, ranging from $+0.0325$ for common words to $+0.875$ for single-3. Family differences are retrieval $+0.64453$, tracking $+0.37000$, aggregation approximately $+0.16208$, and QA approximately $+0.27563$.

Neither arm has empty outputs. TailSpline and MrPro terminate with EOS on $257$ and $248$ of $390$ inputs, respectively, and hit the cap on $133$ and $142$ ($34.1\%$ and $36.4\%$). PPL uses the same document sources as Llama with OLMo-tokenized prefixes. Report and original stream identities accompany the portable aggregates.
\FloatBarrier

\begin{figure}[ht]
\centering\includegraphics[width=\linewidth]{fig_tailspline_details.pdf}
\caption{\textbf{NIAH and language-modeling responses.} NIAH is a subset of each model's depth-balanced RULER-13 panel. Both columns show TailSpline relative to MrPro; PPL changes use each length's absolute PPL. Llama's $8/16$K PPL slightly increases, while OLMo improves at every measured length. The panels have separate vertical scales.}
\label{fig:tailspline-details}
\end{figure}
\FloatBarrier

\subsection{Evaluation implementation}
\label{sec:tailspline-implementation}
Generation is greedy, with repetition penalty $1$, no repeat-ngram restriction and the model's EOS IDs. Llama uses batch size one with cached $8192$-token prefill chunks; OLMo uses batches of at most four identical-length inputs. The static FP32 frequency table is installed after model loading and its installed values are checked against the receipt. NIAH point-depth tasks use two prompts at each target depth $10/30/50/70/90\%$ per length. Other tasks keep their recorded task-specific output budgets.

\paragraph{Input construction.}
The generator pins NVIDIA RULER revision \texttt{c3f5e3b}. Each NIAH task--length cell selects $10$ rows from $64$ candidates using only evidence positions and row identifiers. For single-evidence tasks, the five target depths above each receive two prompts; selection minimizes absolute depth error with tolerance $0.10$. Multi-evidence tasks alternate dispersed targets $(0.10,0.35,0.65,0.90)$ and middle-cluster targets $(0.35,0.45,0.55,0.65)$, with mean absolute tolerance $0.18$. Depths use token positions divided by the context cap, after padding. Task scores are not used in selection.

The filler unit is \texttt{Irrelevant archive padding record; it does not answer the question.} followed by a newline, tokenized without special tokens. Its token sequence is repeated and truncated to the required padding count, then inserted immediately after the user header, before the original task template, until the input occupies cap minus $128$ tokens. The original context, question and references are preserved. In the pinned aggregation templates, the numbered list or coded text follows the task instruction; filler is placed before that instruction. The task-specific generation budgets remain $128$ (NIAH), $30$ (tracking), $120$ (common words), $50$ (frequent words) and $32$ (QA). The row records store padding-token counts and achieved depth errors.

The Llama and OLMo preparation seeds are $20260924$ and $20260925$; each cell adds $300000+10000i+100j$ for length index $i$ and task index $j$. QA uses question-index blocks starting at $5000$, $5010$ and $5020$ for the three lengths. These rules specify sample selection within each model; the two model panels have separate tokenizations. Both exact tables were fixed before their corresponding comparison; historical Llama and OLMo experiments informed the design direction.

\paragraph{Natural-text sources.}
Proof-Pile uses the original \texttt{hoskinson-center/proof-pile} test archive, selecting the first $32$ arXiv records with at least $300000$ characters and retaining up to the first million characters. PG-19 uses individual test-book prefixes. The frozen $46$-document source manifest records the selected $32$ Proof-Pile and $14$ PG-19 identities; each tokenizer takes a prefix of maximum evaluation length plus one token. Documents are scored separately rather than concatenated.

Full-13 intervals use $20{,}000$ paired resamples within each task--length cell, retaining fixed task weights; NIAH intervals aggregate the corresponding retrieval tasks. The report treats different length cells separately for this panel. PPL intervals instead resample the $46$ document identities jointly across all lengths, also using $20{,}000$ draws. The two uncertainty calculations therefore preserve different experimental units. These are conditional evaluation intervals, not uncertainty across training seeds or model families.
